\documentclass[12pt]{article}
\usepackage[utf8]{inputenc}
\usepackage[T1]{fontenc}
\usepackage{hyperref}
\usepackage{amsmath,amssymb,booktabs}
\usepackage{graphicx}
\usepackage{geometry}
\usepackage{listings}
\usepackage{xcolor}
\usepackage{enumitem}
\usepackage{fancyhdr}
\usepackage{float}
\usepackage{caption}
\usepackage{url}

\geometry{margin=1in}
\setlength{\headheight}{14pt}

\hypersetup{
  colorlinks=true,
  linkcolor=blue,
  urlcolor=blue,
  citecolor=blue
}

\definecolor{codebg}{rgb}{0.96,0.96,0.96}
\definecolor{codeframe}{rgb}{0.7,0.7,0.7}
\definecolor{kwcolor}{rgb}{0.0,0.3,0.7}
\definecolor{strcolor}{rgb}{0.6,0.1,0.1}
\definecolor{cmtcolor}{rgb}{0.2,0.5,0.2}

\lstdefinestyle{pycode}{
  language=Python,
  backgroundcolor=\color{codebg},
  frame=single,
  rulecolor=\color{codeframe},
  basicstyle=\footnotesize\ttfamily,
  keywordstyle=\color{kwcolor}\bfseries,
  stringstyle=\color{strcolor},
  commentstyle=\color{cmtcolor}\itshape,
  numbers=left,
  numberstyle=\tiny\color{gray},
  numbersep=6pt,
  breaklines=true,
  breakatwhitespace=false,
  showstringspaces=false,
  tabsize=4,
  xleftmargin=14pt,
}

\lstdefinestyle{jsoncode}{
  backgroundcolor=\color{codebg},
  frame=single,
  rulecolor=\color{codeframe},
  basicstyle=\footnotesize\ttfamily,
  breaklines=true,
  showstringspaces=false,
  tabsize=2,
  xleftmargin=14pt,
}

\lstdefinestyle{textcode}{
  backgroundcolor=\color{codebg},
  frame=single,
  rulecolor=\color{codeframe},
  basicstyle=\footnotesize\ttfamily,
  breaklines=true,
  showstringspaces=false,
  tabsize=2,
  xleftmargin=14pt,
}

\pagestyle{fancy}
\fancyhf{}
\lhead{\small ECE 498 BH --- LLM Reasoning for Engineering}
\rhead{\small Assignment 2}
\cfoot{\thepage}

\begin{document}

% title block matching assignment template
\begin{center}
  {\LARGE\textbf{ECE 498 BH --- LLM Reasoning for Engineering}}\\[6pt]
  {\large\textbf{Assignment 2}}\\[4pt]
  {\normalsize Due 11:59\,pm, 03/31/2026}\\[10pt]
  {\large Prachit Gupta \quad \texttt{prachit2@illinois.edu}}
\end{center}
\thispagestyle{fancy}

\bigskip
\noindent\textbf{Objective}

\noindent
In this assignment you will choose a non-trivial engineering problem and build an LLM API
pipeline to solve it. Rather than interacting with LLMs through web interfaces, you will write
Python code that calls LLM APIs programmatically and integrates a simulation-based automatic
verifier that checks whether the LLM's solution meets all specifications---without any manual
inspection. This assignment trains you in the full loop of LLM-for-engineering: problem
formulation, automated solution generation, and automated evaluation.

\medskip
\noindent\textbf{Submission:} Email me your solutions, including all code files listed in the
Submission Checklist.

% =============================================================================
\newpage
\section{Area Selection [5 pts]}
% =============================================================================

\noindent\textbf{Instructions:} Choose a field that you are familiar with and specify it clearly
below. Examples include (but are not limited to): power electronics, circuit design, control
engineering, signal processing, computer architecture, FPGA design, embedded systems,
aerodynamics, fluid mechanics, structural mechanics, robotics, thermal dynamics, mechatronics,
vibration analysis, manufacturing, and transportation systems.

\bigskip
\noindent\textbf{Selected Engineering Field:}

\medskip
\noindent
Robotics --- Motion planning for Unmanned Aerial Vehicles (UAVs).

% =============================================================================
\newpage
\section{Problem Definition [5 pts]}
% =============================================================================

\noindent\textbf{Instructions:} Provide a well-defined, concrete engineering problem. Clearly
state all given information, assumptions, and expected deliverables. You are not allowed to use
any LLMs for creating this problem. You are allowed to reuse the problem from Assignment 1.

\medskip
\noindent\textit{Calibration hint: Choose a problem where you would expect an LLM to succeed
roughly 40--80\% of the time on a single attempt. If the problem is too easy, every trial will
pass and Sections 6--7 will be uninteresting. If it is too hard, you will have no successful
outputs to analyze.}

\bigskip
\noindent\textbf{Problem Statement:}

\medskip
\noindent\textbf{Note:} This problem is a subsection of a larger research problem I am currently
working on. The entire problem formulation, including all relevant code and the ground truth MPC
reference trajectory, is available at:
\begin{center}
  \url{https://github.com/prachitgupta/vlm-conformal.git}
\end{center}

\subsection*{Local Waypoint Planning for an Autonomous Quadrotor Using Depth Sensing and
Odometry}

\noindent\textbf{Problem Overview.}
LLMs are prompted to solve a high-level motion-planning problem for an autonomous quadrotor that
must reach a goal waypoint while avoiding obstacles \emph{without access to a global map}. The
planner (here LLM) receives real-time depth camera data and vehicle odometry, and is asked to
generate a short sequence (here 5) of \textbf{collision-safe local waypoints} in NED coordinates
that can be tracked by a lower-level off-the-shelf autopilot.

Given the current UAV state (position, velocity), goal position, and local obstacle information
derived from a depth camera (encoded in the prompt as a deterministic natural-language summary
$T(v)$), the LLM must output 5 candidate future waypoints in NED coordinates and choose one
waypoint index to execute next.

\medskip
\noindent\textbf{Deliverable from LLM:} A structured object with exactly three fields:
\begin{itemize}
  \item \texttt{waypoints} --- list of 5 waypoint objects \texttt{\{x, y\}} in NED (North East and Down)  metres;
  \item \texttt{selected\_waypoint\_index} ---  float numerical 2d coordinates with following constraints:-
  \begin{itemize}
      \item First waypoint must be local (At most 3m from current pose)
      \item Consecutive waypoints have to be within 0.2m -1 m range 
  \end{itemize}
  \item \texttt{reasoning} --- a brief free-text explanation of the chosen waypoint.
\end{itemize}

\noindent\textbf{Given Information:}
\begin{itemize}
  \item \textbf{Vehicle platform:} Quadrotor with a front-facing depth camera.
  \item \textbf{Available sensing:}
    \begin{itemize}
      \item Depth camera feed (distance to closest obstacle; local obstacle map).
      \item Real-time visual odometry (drone position and velocity).
    \end{itemize}
  \item \textbf{Available state information:}
    \begin{itemize}
      \item Current pose (position and orientation) in a known local/world NED frame.
      \item Current linear velocity.
      \item Goal pose (or goal waypoint) in the same reference frame.
    \end{itemize}
  \item \textbf{Obstacle information extracted from the sensing pipeline:}
    \begin{itemize}
      \item Depth-frame obstacle points.
      \item Locally generated obstacle map (FIFO buffer of NED points from depth
        back-projection).
      \item Nearest obstacle distance and position (NED).
      \item Sector-wise obstacle clearance/distances (far-left, left, center, right, far-right).
    \end{itemize}
\end{itemize}

\noindent\textbf{Assumptions:}
\begin{itemize}
  \item No global map is available; planning is purely local and reactive
    (receding-horizon style).
  \item The depth camera and odometry are time-synchronised sufficiently for local planning.
  \item Coordinate transforms (camera $\to$ body $\to$ NED) are known and calibrated.
  \item Obstacle detections from the depth camera are reliable within the sensor's valid range.
  \item The planner runs repeatedly in real time, updating the waypoint sequence as new sensor
    data arrives.
  \item A lower-level autopilot controller (PX4 position-tracking mode) exists and can track
    the generated local waypoints.
\end{itemize}

% =============================================================================
\newpage
\section{Specifications and Success Criteria [15 pts]}
% =============================================================================

\noindent\textbf{Instructions:} Carefully enumerate the specifications for your problem so that
there is a clear, quantitative criterion for success for each one. There are three types of
success criteria you may use:
\begin{enumerate}
  \item \textbf{Specification-driven:} The problem has a set of specifications that can be
    verified directly by a computer program. For example, metrics such as phase margin or
    settling time can be computed by a Python script.
  \item \textbf{Testing-driven:} Success is verified by running several predefined tests. For
    instance, if the task is to write code for an engineering application, success can be judged
    by whether the code executes correctly and passes a suite of prescribed test cases.
  \item \textbf{LLM-as-judge:} When no clear-cut quantitative criterion exists, multiple LLMs
    can judge whether the solution is reasonable, and a voting scheme produces a final score.
\end{enumerate}
For engineering problems, we strongly recommend the first two types. Your task is to identify
such criteria for your problem and list them below.

\bigskip
\noindent\textbf{Specifications and Criteria:}

\medskip
The success criteria are \textbf{specification-driven}: every metric is computable by a Python
script from the structured LLM output and the environment vector without running PX4 or Gazebo.
All four criteria below must pass for an overall PASS.

\paragraph{Criterion 1 --- Format / schema compliance (mandatory).}
\begin{itemize}
  \item \texttt{waypoints} is a list of exactly 5 objects, each with finite float fields
    \texttt{x}, \texttt{y}.
  \item \texttt{selected\_waypoint\_index} $\in [0, 4]$.
  \item \texttt{reasoning} is a non-empty string.
\end{itemize}
This criterion is trivially satisfied when using \texttt{instructor}-enforced structured output;
it remains as a guard against schema-level value errors (e.g.\ \texttt{NaN}, out-of-range index).

\paragraph{Criterion 2 --- Kinematic plausibility (hard constraint).}
Let $p_0$ be the current position and $p_1,\ldots,p_5$ the five waypoints. With inter-waypoint
time $\Delta t = 1\,\text{s}$:
\[
  \ell_k = \|p_k - p_{k-1}\|,\quad
  \hat{v}_k = \frac{\ell_k}{\Delta t},\quad
  \hat{a}_k = \frac{|\hat{v}_k - \hat{v}_{k-1}|}{\Delta t}
\]
\textbf{Pass} iff $\hat{v}_k \leq v_{\max} = 5\,\text{m/s}$ and
$\hat{a}_k \leq a_{\max} = 3\,\text{m/s}^2$ for all $k$, and the first leg satisfies
$\ell_1 \leq 3.0\,\text{m}$ (locality constraint).

\paragraph{Criterion 3 --- Goal progress (hard constraint).}
\[
  \Delta d_{\text{sel}} = \|p_0 - g\| - \|p_{\text{sel}} - g\|,\quad
  \Delta d_5 = \|p_0 - g\| - \|p_5 - g\|
\]
\textbf{Pass} iff $\Delta d_{\text{sel}} > 0$ \emph{and} $\Delta d_5 > 0$ (both the selected
waypoint and the final waypoint must strictly reduce distance to goal).

\paragraph{Criterion 4 --- Obstacle safety (core safety metric).}
For each waypoint $p_k$, compute the minimum Euclidean distance to any point in the local
obstacle snapshot $\mathcal{S}_{\text{obs}}$:
\[
  c_k = \min_{\mathbf{o}\in\mathcal{S}_{\text{obs}}} \|p_k - \mathbf{o}\|
\]
\textbf{Pass} iff $\min_k c_k \geq r_{\text{safe}} = 0.8\,\text{m}$.

\medskip
\noindent\textit{Quality metrics reported but not affecting pass/fail:} smoothness
$\kappa = \sum_{k=2}^{5}\|p_{k+1} - 2p_k + p_{k-1}\|$, monotone distance reduction across all
5 waypoints, and selected-index composite score (safety $\times$ progress).

% =============================================================================
\newpage
\section{Simulation-Based Automatic Verifier [35 pts]}
% =============================================================================

\noindent\textbf{Instructions:} Implement a Python function
\texttt{verify(solution) -> dict} that:
\begin{enumerate}
  \item Receives a structured solution object. (You will define the Pydantic schema in
    Section~5; for now, assume the solution arrives as a typed Python object---no string
    parsing required.)
  \item Checks each specification from Section~3 programmatically and returns a structured
    result indicating Pass/Fail per criterion and an overall Pass/Fail.
\end{enumerate}
Your verifier must evaluate the success criteria defined in Section~3 without any human
intervention. The return value should be a dictionary with at minimum the keys \texttt{"pass"}
(bool), \texttt{"details"} (dict of per-specification results), and \texttt{"reason"} (str
summary). The \texttt{"details"} and \texttt{"reason"} fields are not just for your own
debugging---in Section~7 you will feed \texttt{"reason"} back to the LLM so it knows what
failed and by how much, enabling meaningful self-correction.

\medskip
\noindent\textit{Important: You are allowed to use LLMs to help you develop the verifier code,
but you must manually verify that the code is correct. In addition, you must run it on at least
two hand-crafted test cases---one passing and one failing---and include those test results below
to confirm correctness.}

\bigskip
\noindent\textbf{Verifier Code:} (attach as \texttt{verifier.py})

\lstinputlisting[style=pycode, caption={\texttt{verifier.py}}]{verifier.py}

\bigskip
\noindent\textbf{Manual Verification Test Results:}

\medskip
\noindent\textbullet\quad\textbf{Test Case 1 (expected: PASS):}

\medskip
Five forward-progressing waypoints starting at $(2.7,\,0.35,\,-2.5)$, which is $2.85\,\text{m}$
from $p_0 = (-0.14,\,0.31,\,-2.43)$ --- within the $3.0\,\text{m}$ locality limit.
Max implied speed $2.62\,\text{m/s} \leq 5\,\text{m/s}$, max accel
$1.35\,\text{m/s}^2 \leq 3\,\text{m/s}^2$, min obstacle clearance $4.02\,\text{m} \geq 0.8\,\text{m}$,
final waypoint $9.28\,\text{m}$ closer to goal. All criteria pass. Verifier output:

\begin{lstlisting}[style=jsoncode]
{
  "pass": true,
  "details": {
    "kinematics":    {"pass": true,
                      "segment_lengths_m": [2.85, 2.53, 1.95, 2.25, 2.52]},
    "goal_progress": {"pass": true, "progress_selected_m": 2.71,
                      "progress_final_m": 9.28},
    "safety":        {"pass": true, "min_clearance_m": 4.02},
    "smoothness":    {"second_diff_norm": 0.44, "max_turn_deg": 18.3}
  },
  "reason": "All criteria passed."
}
\end{lstlisting}

\medskip
\noindent\textbullet\quad\textbf{Test Case 2 (expected: FAIL):}

\medskip
First waypoint at $(6.0,\,0.3,\,-2.5)$ is $6.14\,\text{m}$ from $p_0$, violating both the
$3.0\,\text{m}$ locality constraint and $v_{\max} = 5\,\text{m/s}$. Verifier correctly returns
FAIL with a precise, human-readable reason string. Verifier output:

\begin{lstlisting}[style=jsoncode]
{
  "pass": false,
  "details": {
    "kinematics": {
      "pass": false,
      "violations": [
        "First leg 6.14 m > 3.0 m (locality constraint)",
        "Segment 0: speed 6.14 m/s > 5.0 m/s"
      ]
    },
    "goal_progress": {"pass": true},
    "safety":        {"pass": true}
  },
  "reason": "KINEMATICS: First leg 6.14 m > 3.0 m (locality constraint);
             Segment 0: speed 6.14 m/s > 5.0 m/s"
}
\end{lstlisting}

% =============================================================================
\newpage
\section{LLM API Pipeline with Enforced Structured Output [15 pts]}
% =============================================================================

\noindent\textbf{Instructions:} Choose one model from the following APIs: GPT 5.4, Gemini
3.0 Pro, or Claude 4.6 Opus. Write a Python script that:
\begin{enumerate}
  \item Defines a Pydantic model describing the expected solution schema.
  \item Uses the \texttt{instructor} library to patch the API client so the LLM is forced to
    return a valid instance of that Pydantic model rather than free-form text. This completely
    removes the need for fragile string parsing in your verifier.
  \item Passes the structured solution object directly to \texttt{verify()} from Section~4.
  \item Prints the verification result.
\end{enumerate}
\textit{Why instructor? By enforcing a schema at the API call level, you guarantee that
\texttt{verify()} always receives clean, typed values. This pattern is standard practice in
production LLM-for-engineering pipelines.}

\bigskip
\noindent\textbf{Pipeline Code:} (attach as \texttt{pipeline.py})

\medskip
\noindent\textbf{Chosen model:} \texttt{gpt-4o-mini} via the OpenAI API with
\texttt{instructor} structured-output enforcement.

\lstinputlisting[style=pycode, caption={\texttt{pipeline.py}}]{pipeline.py}

% =============================================================================
\newpage
\section{Repeated Sampling [10 pts]}
% =============================================================================

\noindent\textbf{Instructions:} Extend your pipeline to query the chosen LLM five times in
independent API calls (i.e., no shared conversation history). Run the verifier automatically on
each response and record results in the table below.

\medskip
\noindent\textit{Note: All five trials must be run programmatically in a loop---no manual
copy-paste. Submit your looping script as \texttt{repeated\_sampling.py}. Your code should
automatically print the Pass/Fail results for each trial; report those results in the table
below.}

\lstinputlisting[style=pycode, caption={\texttt{repeat\_sampling.py}}]{repeat_sampling.py}

\begin{table}[H]
  \centering
  \begin{tabular}{ccc}
    \toprule
    \textbf{LLM} & \textbf{Trial} & \textbf{Pass/Fail} \\
    \midrule
                     & 1 & Pass \\
                     & 2 & Pass \\
    GPT-4o-mini     & 3 & Pass \\
                     & 4 & Pass \\
                     & 5 & Pass \\
    \bottomrule
  \end{tabular}
  \caption{Repeated sampling results. Pass/Fail determined automatically by the verifier.}
  \label{tab:repeated_sampling}
\end{table}

\medskip
\noindent
For dataset row~0, all five repeated-sampling trials passed. Each trial produced the same
straight center-progress waypoint sequence: $(0.0, 0.5)$, $(0.0, 1.0)$, $(0.0, 1.5)$,
$(0.0, 2.0)$, and $(0.0, 2.5)$, with only small wording changes in the reasoning field. This
shows that repeated independent calls were highly stable for the clear-corridor case.

% =============================================================================
\newpage
\section{Self-Refinement Loop [5 pts]}
% =============================================================================

\noindent\textbf{Instructions:} Implement a multi-turn API conversation where the LLM
iteratively refines its own solution based on verifier feedback. Use \texttt{instructor}
throughout so every turn returns a typed Pydantic object (e.g., your \texttt{Solution} model
from Section~5). Specifically:
\begin{enumerate}
  \item In turn 1, ask the LLM to solve the problem (structured output enforced).
  \item Run the verifier on the returned solution object.
  \item If the solution fails, append the verifier's \texttt{"reason"} string to the conversation
    history and ask the LLM to fix the solution. Repeat for up to 2 refinement turns.
  \item Record the Pass/Fail status after each turn.
\end{enumerate}
Run this loop once. Submit your script as \texttt{self\_refinement.py}. The verifier feedback
passed back to the LLM must come directly from your \texttt{verify()} function---it must not be
written by hand.

\lstinputlisting[style=pycode, caption={\texttt{self\_refinement.py}}]{self_refinement.py}

\begin{table}[H]
  \centering
  \begin{tabular}{ccccc}
    \toprule
    \textbf{LLM} & \textbf{Turn 1} & \textbf{Turn 2} & \textbf{Turn 3 (if needed)} &
    \textbf{Final Pass/Fail} \\
    \midrule
    GPT-4o-mini & Fail & Fail & Fail & \textbf{Fail} \\
    \bottomrule
  \end{tabular}
  \caption{Self-refinement results. All Pass/Fail entries are determined by the automatic
  verifier.}
  \label{tab:self_refinement}
\end{table}

\medskip
\noindent\textbf{Brief Observations (1--3 sentences):} For dataset row~4, self-refinement did
not recover a passing solution. All three turns returned the same waypoint sequence
$(0.68, 2.06)$ through $(0.68, 4.06)$, and every turn failed the same safety check with minimum
nearest-obstacle-proxy clearance $0.51\,\text{m} < 0.80\,\text{m}$. In this case, passing the
verifier feedback back to the model was not sufficient to change the trajectory.

% =============================================================================
\newpage
\section{Final Analysis and Insights [10 pts]}
% =============================================================================

\noindent\textbf{Instructions:} Discuss any additional findings based on your experiments.
Consider the following questions:
\begin{itemize}
  \item What were the main failure modes of the LLM (e.g., wrong parameter values, unstable
    designs)? How did the verifier's feedback messages help or hinder refinement?
  \item How reliable was the LLM across trials?
  \item Did enforcing structured output via \texttt{instructor} eliminate parsing errors
    entirely, or did you still encounter issues (e.g., values out of range, schema validation
    failures)? Describe any edge cases.
  \item Suggest one concrete improvement to either the Pydantic schema design, the prompt, or
    the verifier feedback format that you believe would increase the pass rate. Justify your
    answer.
  \item (Optional) If you tried more advanced prompting strategies (chain-of-thought, few-shot
    examples, richer \texttt{Field(description=...)} annotations), describe their effect.
\end{itemize}

\bigskip
\noindent\textbf{Main failure modes.}
In the recorded OpenAI runs, the dominant failure mode was the safety criterion on the
\texttt{row-index 4} caution scenario, not kinematics. The model repeatedly produced the
waypoints $(0.68, 2.06)$, $(0.68, 2.56)$, $(0.68, 3.06)$, $(0.68, 3.56)$, and $(0.68, 4.06)$,
which satisfied locality, speed, acceleration, and goal-progress constraints but approached the
nearest-obstacle proxy too closely. The verifier reported a minimum clearance of
$0.512\,\text{m}$, below the required $0.8\,\text{m}$.

\medskip
\noindent\textbf{Reliability across trials.}
GPT-4o-mini achieved a 5/5 (100\%) first-attempt pass rate on repeated sampling for
\texttt{row-index 0}. Across all five independent calls, the waypoint coordinates were identical
and only the free-text reasoning varied slightly. This indicates that for a clear corridor prompt
the model is highly deterministic at the chosen settings.

\medskip
\noindent\textbf{Structured output via \texttt{instructor}.}
Enforcing the Pydantic schema via \texttt{instructor} completely eliminated format and
JSON-parsing errors. Across the recorded calls used here (two single pipeline runs, five
repeated-sampling calls, and three self-refinement turns), there were zero schema-validation
failures: every response contained exactly 5
finite-valued waypoints and a valid index. This is a significant improvement over the HW1
web-UI workflow, where markdown fences and missing keys required pre-processing. The only
remaining failures were semantic (constraint violations), which is exactly what the verifier is
designed to catch. No edge cases involving \texttt{NaN} or out-of-range indices were observed,
though the \texttt{\_check\_format} guard remains in the verifier as insurance.

\medskip
\noindent\textbf{Suggested improvement.}
Adding a \texttt{scratchpad: str} field to the Pydantic model and instructing the model to
``first compute all segment lengths and implied speeds in the scratchpad, then populate
\texttt{waypoints}'' would front-load constraint verification inside the generation step. This
directly targets the observed kinematic failure: the model would catch its own speed violation
before committing to coordinates (chain-of-thought with constraint pre-computation). The
\texttt{scratchpad} field would be excluded from the verifier, adding no downstream complexity.

\medskip
\noindent\textbf{(Optional) Advanced prompting.}
The few-shot demonstrations in the system prompt (Examples A, B, C) were carried over verbatim
from HW1 and contributed to the consistent output format across all trials. Richer
\texttt{Field(description=...)} annotations on the \texttt{Waypoint} fields (e.g.,
\textit{``North coordinate; ensure $\|p_1 - p_0\| \leq 3.0\,\text{m}$ before placing this
waypoint''}) were trialled during development and reduced first-attempt kinematic violations
in informal testing at no additional runtime cost.

% =============================================================================
\newpage
\subsection*{Submission Checklist}
% =============================================================================

Please confirm that your submission includes all of the following:

\begin{itemize}[label=$\boxtimes$]
  \item \textbf{SETUP.md} --- environment and dependency documentation. Python~3.11+;
    packages: \texttt{openai$\,\geq\,$2.0.0}, \texttt{instructor$\,\geq\,$1.4.0},
    \texttt{pydantic$\,\geq\,$2.7.0}. No external simulation
    tools required for Layer~A verification.
  \item The completed PDF/LaTeX files with all tables and written responses filled in.
  \item \textbf{verifier.py} --- standalone verifier with manual test results (Section~4).
  \item \textbf{pipeline.py} --- single-query pipeline with \texttt{instructor} structured
    output (Section~5).
  \item \textbf{repeated\_sampling.py} --- five-trial loop (Section~6).
  \item \textbf{self\_refinement.py} --- multi-turn refinement loop with structured output
    (Section~7).
\end{itemize}

\noindent API keys do not appear in any submitted file; they are loaded via
\texttt{os.environ["OPENAI\_API\_KEY"]}.

% =============================================================================
\newpage
\section*{Appendix: Added Materials}
% =============================================================================

\noindent\textbf{Fixed Prompt Used for Assignment 2}

\medskip
\noindent This is the fixed planning-policy prompt used by the current codebase. It is loaded by
\texttt{pipeline.py} from \texttt{assignment2/fixed\_prompt.txt} and stays unchanged across all
trials. It defines the task, the 2D waypoint schema, the waypoint-spacing constraints, and the
few-shot demonstration.

\begin{lstlisting}[style=textcode,caption={assignment2 fixed prompt used for the API pipeline}]
You are the motion planner for an autonomous quadrotor.

SECTION 1: CONTEXT AND INSTRUCTIONS
Task:
- Plan a short local trajectory of the next 5 waypoints.
- The 5-waypoint sequence should be consistent, smooth, collision-safe, and useful as a local trajectory label.
- Use all provided context, including odometry-derived state and obstacle/depth cues.
- Plan only in the horizontal plane: choose waypoint x and y values.
- Do not output z. The planner will add z from the goal altitude after parsing your response.

Frame and units:
- Coordinate frame is NED (North-East-Down), meters.
- x: North, y: East.
- Output waypoint objects must contain only x and y fields.

Operational constraints:
- Hard safety: never place waypoints inside or too close to obstacles.
- Keep at least 1.5 m obstacle clearance when feasible.
- Keep waypoint speeds at or below 5 m/s and waypoint-to-waypoint acceleration changes at or below 3 m/s^2.
- Keep waypoints dynamically feasible and smooth (no abrupt reversals or zig-zags).
- Ensure forward progress toward the goal.
- First waypoint MUST BE LOCAL: AT MOST 3 m from current position.
- Consecutive waypoint spacing should usually be 0.5 m to 3.0 m unless safety requires shorter moves.
- Return the index of the waypoint to execute next.

Planning objective priority:
1) Safety and feasibility
2) Progress to goal
3) Smoothness and efficiency
4) Consistent 5-point local trajectory output

SECTION 2: FEW-SHOT DEMONSTRATION
Example A (clear corridor):
Input summary:
- current_position: (-0.02, 0.01, -2.16)
- goal_position:  (0.00, 29.00, -2.50)
- obstacle cue: Nearest obstacle sector is center at 4.02 m and clearance status is clear.

Interpretation:
- The center corridor is clear and supports direct forward progress.
- The waypoints should stay local, smooth, and advance along the goal direction.
- selected_waypoint_index should indicate the immediate next step.

Output:
{
  "waypoints": [
    {"x": 0.19, "y": 0.44},
    {"x": 0.39, "y": 0.89},
    {"x": 0.58, "y": 1.33},
    {"x": 0.78, "y": 1.78},
    {"x": 0.97, "y": 2.22}
  ],
  "selected_waypoint_index": 0,
  "reasoning": "Center corridor is clear, so a smooth forward trajectory is safe."
}

SECTION 3: EXPECTED OUTPUT
Return ONLY one valid JSON object. No markdown. No extra text.

Use this Pydantic-style object:
```python
from pydantic import BaseModel

class Waypoint(BaseModel):
    x: float
    y: float

class WaypointSolution(BaseModel):
    waypoints: list[Waypoint]
    selected_waypoint_index: int
    reasoning: str
```

Output rules:
- "waypoints" must contain exactly 5 waypoint objects.
- Every waypoint must include numeric x and y.
- Waypoints must be ordered in execution order from immediate next step to farther local lookahead.
- The first waypoint must be within 3.0 m of the current position.
- Consecutive waypoint spacing should usually be 0.5 m to 3.0 m unless safety requires shorter moves.
- "selected_waypoint_index" must be an integer from 0 to 4.
- "reasoning" should briefly explain why the selected waypoint is safe and useful.
\end{lstlisting}

\bigskip
\noindent\textbf{Environment Prompt Used in Verifier Test Example}

\medskip
\noindent This prompt is generated from dataset row~0 using the same
\texttt{load\_environment\_prompt(...)} helper that the pipeline uses. It is the clear-corridor
case used for the successful single pipeline run and for all five repeated-sampling trials.

\begin{lstlisting}[style=textcode,caption={Dataset prompt used in the local verifier example}]
Environment section (T(v)):
- Current position NED is (-0.02, 0.01, -2.16) m.
- Current velocity is (-0.00, -0.00, -0.00) m/s (speed 0.00 m/s).
- Goal position NED is (0.00, 29.00, -2.50) m.
- Output contract: predict exactly 5 waypoint x/y pairs in NED. The planner will set every waypoint z to goal z = -2.50 m before execution.
- Goal delta from current state is (0.02, 28.99, -0.34) m with distance 28.99 m.
- Obstacle pipeline: depth_camera_to_body_to_ned_current_frame.
- Current depth frame obstacle points projected into NED: 11. Nearest depth obstacle in NED is at (-0.36, 4.15, -2.13) m, relative vector (-0.35, 4.15, 0.03) m, distance 4.16 m.
- Depth validity fraction is 0.61. Global nearest obstacle distance is 4.02 m.
- Sector minimum distances [far_left, left, center, right, far_right] are [4.37, 4.37, 4.02, 4.38, 4.38] m.
- Nearest obstacle sector is center at 4.02 m and clearance status is clear.
- Lateral planning hint: prefer center progress.
\end{lstlisting}

\bigskip
\noindent\textbf{Additional Environment Prompt Collected for API Runs}

\medskip
\noindent This prompt is generated from dataset row~4 using the same pipeline loader. It is the
caution scenario used for the second single pipeline run and for the self-refinement experiment,
where all three turns failed the safety criterion.

\lstinputlisting[style=textcode,caption={Second dataset prompt selected for assignment2 experiments}]{generated/example_env_row4.txt}

\bigskip
\noindent\textbf{Updated Code Used in This Assignment 2 Workspace}

\lstinputlisting[style=pycode,caption={\texttt{assignment2/verifier.py}}]{verifier.py}

\lstinputlisting[style=pycode,caption={\texttt{assignment2/pipeline.py}}]{pipeline.py}

\lstinputlisting[style=pycode,caption={\texttt{assignment2/repeat\_sampling.py}}]{repeat_sampling.py}

\lstinputlisting[style=pycode,caption={\texttt{assignment2/self\_refinement.py}}]{self_refinement.py}

\end{document}
